\epigraph{The brain is a wonderful organ; it starts working the moment you get up in the morning and does not stop until you get into the office}
{\textit{Robert Frost}}

\section{Overview}

% Talk about how lectures are gonna be different to classroom learning and how to deal with that, how the material comes up quite fast. 
% how many lectures per module 
% what is the general structure like of lectures 
% Note on notes being provided beforehand + other resources 
% how long they are 

% timetable can be found on Seats and on blackboard - useful to link timetable to google calender or something equivalent. 
% Check in with Seats - QR code or the seats code will usually be given start of the lecture  - keeps a good track of attendance 



You'll become familiar with two names very quickly: \textbf{Blackboard} and \textbf{Seats}. 

Blackboard is the university's online learning platform where lectures notes are uploaded, announcements are posted and the coursework and recordings are stored. If something official happens in the module, it will most likely appear there. 

One of the most useful things you can do in the first week is link your timetable to your own calender, whether that's Google calender or something equivalent. Remember that universities will not chase you the way school did. There are no bells and no-one reminds you every morning where you're supposed to me, so the responsibility shifts to you. 

Seats is the system used for timetabling and attendance. It shows your lecture locations, times, and room numbers. It's also how attendance is recorded. At the start of lectures, a QR or Seats code will be displayed on the board which students can scan or type to register their attendance on the app. 

It is worth getting into the habit of checking both early in the semester because you might find it may unnecessary stress. Knowing where things are in advance means you can focus on the maths rather than the logistics. 

\subsection{Core Skills}

% info on core skills module? 

\subsection{Problem Classes}
% What they are - how to make the most of them 
% Try and have a look at the problems beforehand 
% what is the mode up-to-date info on this? is this somethign that continues after the 1st yr or not? 

\section{Recorded Doesn't Mean Optional}

You will quickly realise that most lectures are recorded, and at some point you might be temped to think, `I can just watch it later', especially on a cold and rainy morning, or after a late night, or when the idea of sitting in a lecture feels mildly exhausting. 

I cannot stress this enough, but being physically present in a lecture matters more than you think. When you are in the room, you are less likely to pause a video and get distracted by your phone. You are also less likely to tell yourself that you'll `finish it later' and risk not coming back to it. You are more likely to stay with the explanations as they unfold, even if you don't fully understand them yet. You hear the throwaway comments that never make it into typed notes, and pick up on the emphasis in the lecturer's voice when they say something is important, and that emphasis stays with you. 

There is also something about the collective focus that recordings cannot replicate. Sitting in a room full of people who are all attempting to follow the same proof creates a concentration that watching it alone at 1.5x speed, half distracted cannot compete with. In addition to that, you are able to ask questions during or after the lecture if you attend in person. 

Recordings are useful, in fact, they are incredibly useful for revision or revisiting a proof if that felt unclear the first time. They are a safety net, but not a replacement for showing up. 

That said, there will be weeks where life happens. Missing a lecture isn’t the end of the world. Just don’t let `I’ll watch it later' turn into `I’m now three weeks behind and slightly afraid to look at the video'. The longer you rely on your future self to catch up, the easier it becomes to drift out of sync with the module, and once you're out of sync, everything feels harder than it needs to be. 

So if you can be there, be there. 

\section{When Your Brain Just Isn't Working}

There often comes a point in almost every student's journey where one morning, for no dramatic reason, your brain simply refuses to cooperate. 

The lecturer is speaking, and the words are technically English, and the notation seems familiar. The content might not even be complicated, and yet, nothing seems to land. You realise you've copied half a page without actually processing a single line and catch yourself been thinking about what you're going to have for lunch for the last 8 minutes. 

Rest assured, this is normal and happens to the best of us. It does not mean you have suddenly lost the ability to think mathematically. Engaging with university mathematics requires sustained concentration in a way that school rarely did. The abstraction is higher and the cognitive load is heavier. Your brain will sometimes hit its limit, especially if you're tired, stressed, slightly underslept, or just adjusting to a new routine and environment. 

Not every lecture will feel inspiring and not every topic will immediately excite you. On days when your performance is nowhere near optimal, lower the bar slightly. Instead of aiming to understand everything, narrow the target. Pick one line or sentence on the board and ask yourself what it is actually saying. Or write down one question you would ask if you were brave enough, or underline one step in the proof that feels unclear and mark it to revisit later. 

If your attention in slipping, change your posture. Sit up straight and just listen without feeling the need to copy everything down. Just write a short sentence that summarises what you have learned so far, or questions if you're confused about something. Small physical adjustments often help reset focus more than you'd expect. 

And if you're brain truly isn't cooperating despite your best efforts, just accept that for what it is. Stay present as best you can, then review the material later when you're more alert. Not all the understanding needs to happen in one sitting in the lecture hall. 

The key message here being, one unfocused lecture does not define your trajectory\footnote{Pun intended}. Boredom, distraction, fatigue- it's all part of the learning experience, and learning to work with them rather than against them is part of becoming a good learner. 


\subsection{Catching Up On Lectures}

At some point, you may very well miss a lecture. Maybe you were ill, overslept, or just needed a morning off. It happens, and one missed lecture is not a crisis.

What turns it into a crisis is delay. 

One missed lecture feels manageable, two feel slightly annoying and by week five, you might start avoiding the recordings or looking at notes entirely because it just feels overwhelming. 

The trick to avoid going down that spiral is quite simple: catch up quickly. If you miss a lecture, watch the recording within a day or two while the topic is still current and notation still looks familiar. Don't let it all pile onto future you. If for some reason due to time constraints you can't watch the recordings, ask your coursemates or friends about which point in the notes the lecturer has gone to and read up to that point. 

When you do catch up, don't treat the recording like a background noise. Pause and try to follow the argument yourself. When a definition appears, stop and ask whether you understand what each part is doing and when a proof begins, try and predict the next step before the lecturer reveals it. 

Falling slightly behind is very normal. But letting it snowball is what makes learning more difficult than it needs to be. 


\section{Asking questions}
% Why it's a good thing - link back to being present in person 
% Addressing fear and embarassemtn of asking questions in lectures 

There is often a quiet fear around asking questions in lectures. If you ever find yourself sitting in a lecture room, slightly unsure about a definition or confused about some step in a proof, and you look around the room and everyone else looks calm, and you might think, `Maybe it's just me'. 

It almost never is. Most people are running through the same internal debate of `Is this worth asking? Is this obvious? Is this too stupid of a question?', etc. 

Here's something I've learned: a question does not need to be clever, it just needs to be honest. Questions like `How did we go from Line 1 to Line 2 of this proof', or `Is there an example where this conditions fails?' are perfectly good questions to ask in lecture. There are no stupid questions. 

By asking questions, you not only benefit yourself, but also others in the lecture, including the lecturer themselves. The lecturer can see where the gaps are, other students relax because someone has said what they were thinking, and most importantly, the room becomes less passive. Learning mathematics, doing mathematics, is a process, and asking questions is part of that process. 

If asking during the lecture feels too intimidating, most lecturer's are more than happy for students to ask them at the end of the lecture, or during office hours and workshops. 

\section{Giving feedback}

% Encourage people to give mid semeseter feedback to any feedback like for SSLC when requested- helps people do better 
% Encourage people to do surveys and why they are important 
% give examples on what looks like good feedback vs poor feedback 

At some point during the semester, you will be asked to fill in feedback forms or respond to surveys. 

It is very tempting to either rush through them, or completely ignore them. It is also very easy to write things like `Too fast' or `Confusing' or `Good'. 

But feedback only becomes useful when it is specific. Writing `This lecture is confusing' doesn't give anyone anything to work with, whereas writing `There are too many definitions and very few examples to see how those definitions work' is useful. 

Some responses convey reaction, others information, so ensure that whenever you're in a position of providing feedback, you relay the information over reaction as best as you can. 

Lecturers are not mind readers. They often assume a certain level of clarity because they've been living with the material for far far longer than we have. So if ten students feel quietly lost but say nothing, nothing changes. 

Feedback not only helps the lecturers, but also the cohort. And it helps the individual, because the act of explaining why something was unclear forces you to reflect on what you were expecting. 

\subsection{Too fast? Too Slow?}

There will be lectures where some definitions will arrive in rapid succession and you are still trying to process the first line while the theorem is already on the board. 

A lecture being `fast paced' usually happens when you are encountering abstraction for the first time, which takes time to settle. Exposure of a new topic precedes clarity, and you are not expected to internalise everything immediately. 

On the other hand, there might be lectures that feel slow, perhaps even repetitive, and you may find yourself zoning out because you think you already understand it. That is often the moment where it is worth looking deeper into what it being taught, asking why the definition is phrased in a certain way, or what breaks if an assumption is removed, or even how the idea connects to something you've already seen. 

The pace of a lecture is rarely perfectly aligned with your personal understanding, and that is normal. However, if you feel like something is unreasonable, you should mention that in either one of the feedback forms or by talking to your lecturer. 

\section{Coursework}
% Make a note of when it is due - most modules follow a timetable, eg: setting work every friday and due the next friday 

Most modules will have some form of coursework, which in practice usually means a weekly problem sheet or set of exercises that are released on a fixed day and due the following week. For example, you might get an assignment set on Monday afternoon and due the next Monday, but always check the specific module outline for the details. Some coursework might be on paper, others might be delivered using an online environment like Stack. 

At first, coursework can feel constant. You submit one sheet and the next one appears almost immediately, and it can feel like there is never a complete clear week. Having said that, the regular rhythm of doing problem sheets is actually really helpful as it stops the module from going into the `I'll catch you later' territory and encourages you to engage with the content. 

The key is to treat coursework as part of your weekly structure instead of something extra that appears unexpectedly. Build your week and routine around when the assignments are due, and decide early when you will attempt each one. 

More often than not, the hardest part of coursework is getting started with it. Start working on it earlier than you think. There can be a very common pattern in the first year where the coursework is released on Monday, you glance at it on Wednesday, and realise on Sunday evening that you still haven't attempted exercise 1. 

If you leave assignments until the night before, you reduce the amount of time you have to get stuck productively. The questions are designed to make you think, and thinking takes time. 

So try and get into a routine of starting the coursework at least one day in advance if not before. If you have attempted a question and still feel stuck, the department has built something very useful into the structure of the week that helps with such problems: \textbf{First Year Workshops}. For if and when find yourself on coursework, the next chapter aims to explain how workshops work. 


\subsection{Extensions and Special Considerations}
% Note on applying for special cons + asking for extensions - look up latest info on this - dpeends on the module so better to ask the module lead 


As with everything, you might have moments where multiple deadlines collide in a way that genuinely overwhelms you. Universities know this and that's why extensions and special consideration procedures exist. 

% Add links to uni and dept policy on speical cons for coursework 


